% !TeX root = ../../ONMIPA Alin.tex
Sebelum memasuki pembahasan pada bagian ini, pembaca
diharapkan telah memahami \textit{operasi baris elementer}
(OBE) dan \textit{operasi kolom elementer} (OKE).
Kedua operasi tersebut merupakan dasar penting dalam
menyederhanakan matriks dan menganalisis sifat-sifatnya.

Sebagai pengingat, terdapat tiga jenis operasi baris
elementer, yaitu
\[
    \begin{aligned}
        R_i & \leftrightarrow R_j
            & & \text{(pertukaran dua baris)}, \\
        R_i & \leftarrow \lambda R_i,\quad \lambda\neq 0
            & & \text{(perkalian dengan skalar)}, \\
        R_i & \leftarrow R_i+\lambda R_j,\quad i\neq j
            & & \text{(penjumlahan kelipatan baris)}.
    \end{aligned}
\]
Operasi kolom elementer didefinisikan secara analog.
Setiap OBE dapat direpresentasikan sebagai perkalian
matriks elementer dari kiri, sedangkan OKE sebagai
perkalian matriks elementer dari kanan, yakni
\[
    A\xrightarrow{\mathrm{OBE}} E_iA
    \qquad\text{dan}\qquad
    A\xrightarrow{\mathrm{OKE}} AF_j,
\]
dengan $E_i$ dan $F_j$ merupakan matriks elementer.
